% Options for packages loaded elsewhere
\PassOptionsToPackage{unicode}{hyperref}
\PassOptionsToPackage{hyphens}{url}
\PassOptionsToPackage{dvipsnames,svgnames,x11names}{xcolor}
\documentclass[
  spanish,
  10pt,
  a4paper,
]{article}
\usepackage{xcolor}
\usepackage[top=2.2cm,bottom=2.2cm,left=2.1cm,right=2.1cm]{geometry}
\usepackage{amsmath,amssymb}
\setcounter{secnumdepth}{5}
\usepackage{iftex}
\ifPDFTeX
  \usepackage[T1]{fontenc}
  \usepackage[utf8]{inputenc}
  \usepackage{textcomp} % provide euro and other symbols
\else % if luatex or xetex
  \usepackage{unicode-math} % this also loads fontspec
  \defaultfontfeatures{Scale=MatchLowercase}
  \defaultfontfeatures[\rmfamily]{Ligatures=TeX,Scale=1}
\fi
\usepackage{lmodern}
\ifPDFTeX\else
  % xetex/luatex font selection
\fi
% Use upquote if available, for straight quotes in verbatim environments
\IfFileExists{upquote.sty}{\usepackage{upquote}}{}
\IfFileExists{microtype.sty}{% use microtype if available
  \usepackage[]{microtype}
  \UseMicrotypeSet[protrusion]{basicmath} % disable protrusion for tt fonts
}{}
\makeatletter
\@ifundefined{KOMAClassName}{% if non-KOMA class
  \IfFileExists{parskip.sty}{%
    \usepackage{parskip}
  }{% else
    \setlength{\parindent}{0pt}
    \setlength{\parskip}{6pt plus 2pt minus 1pt}}
}{% if KOMA class
  \KOMAoptions{parskip=half}}
\makeatother
\usepackage{color}
\usepackage{fancyvrb}
\newcommand{\VerbBar}{|}
\newcommand{\VERB}{\Verb[commandchars=\\\{\}]}
\DefineVerbatimEnvironment{Highlighting}{Verbatim}{commandchars=\\\{\}}
% Add ',fontsize=\small' for more characters per line
\newenvironment{Shaded}{}{}
\newcommand{\AlertTok}[1]{\textcolor[rgb]{1.00,0.00,0.00}{\textbf{#1}}}
\newcommand{\AnnotationTok}[1]{\textcolor[rgb]{0.38,0.63,0.69}{\textbf{\textit{#1}}}}
\newcommand{\AttributeTok}[1]{\textcolor[rgb]{0.49,0.56,0.16}{#1}}
\newcommand{\BaseNTok}[1]{\textcolor[rgb]{0.25,0.63,0.44}{#1}}
\newcommand{\BuiltInTok}[1]{\textcolor[rgb]{0.00,0.50,0.00}{#1}}
\newcommand{\CharTok}[1]{\textcolor[rgb]{0.25,0.44,0.63}{#1}}
\newcommand{\CommentTok}[1]{\textcolor[rgb]{0.38,0.63,0.69}{\textit{#1}}}
\newcommand{\CommentVarTok}[1]{\textcolor[rgb]{0.38,0.63,0.69}{\textbf{\textit{#1}}}}
\newcommand{\ConstantTok}[1]{\textcolor[rgb]{0.53,0.00,0.00}{#1}}
\newcommand{\ControlFlowTok}[1]{\textcolor[rgb]{0.00,0.44,0.13}{\textbf{#1}}}
\newcommand{\DataTypeTok}[1]{\textcolor[rgb]{0.56,0.13,0.00}{#1}}
\newcommand{\DecValTok}[1]{\textcolor[rgb]{0.25,0.63,0.44}{#1}}
\newcommand{\DocumentationTok}[1]{\textcolor[rgb]{0.73,0.13,0.13}{\textit{#1}}}
\newcommand{\ErrorTok}[1]{\textcolor[rgb]{1.00,0.00,0.00}{\textbf{#1}}}
\newcommand{\ExtensionTok}[1]{#1}
\newcommand{\FloatTok}[1]{\textcolor[rgb]{0.25,0.63,0.44}{#1}}
\newcommand{\FunctionTok}[1]{\textcolor[rgb]{0.02,0.16,0.49}{#1}}
\newcommand{\ImportTok}[1]{\textcolor[rgb]{0.00,0.50,0.00}{\textbf{#1}}}
\newcommand{\InformationTok}[1]{\textcolor[rgb]{0.38,0.63,0.69}{\textbf{\textit{#1}}}}
\newcommand{\KeywordTok}[1]{\textcolor[rgb]{0.00,0.44,0.13}{\textbf{#1}}}
\newcommand{\NormalTok}[1]{#1}
\newcommand{\OperatorTok}[1]{\textcolor[rgb]{0.40,0.40,0.40}{#1}}
\newcommand{\OtherTok}[1]{\textcolor[rgb]{0.00,0.44,0.13}{#1}}
\newcommand{\PreprocessorTok}[1]{\textcolor[rgb]{0.74,0.48,0.00}{#1}}
\newcommand{\RegionMarkerTok}[1]{#1}
\newcommand{\SpecialCharTok}[1]{\textcolor[rgb]{0.25,0.44,0.63}{#1}}
\newcommand{\SpecialStringTok}[1]{\textcolor[rgb]{0.73,0.40,0.53}{#1}}
\newcommand{\StringTok}[1]{\textcolor[rgb]{0.25,0.44,0.63}{#1}}
\newcommand{\VariableTok}[1]{\textcolor[rgb]{0.10,0.09,0.49}{#1}}
\newcommand{\VerbatimStringTok}[1]{\textcolor[rgb]{0.25,0.44,0.63}{#1}}
\newcommand{\WarningTok}[1]{\textcolor[rgb]{0.38,0.63,0.69}{\textbf{\textit{#1}}}}
\ifLuaTeX
\usepackage[bidi=basic,shorthands=off]{babel}
\else
\usepackage[bidi=default,shorthands=off]{babel}
\fi
\ifLuaTeX
  \usepackage{selnolig} % disable illegal ligatures
\fi
\setlength{\emergencystretch}{3em} % prevent overfull lines
\providecommand{\tightlist}{%
  \setlength{\itemsep}{0pt}\setlength{\parskip}{0pt}}
% Shared visual layer for the generated Kaleido FlowTwin reports.
\usepackage{helvet}
\renewcommand{\familydefault}{\sfdefault}
\usepackage{microtype}
\usepackage{xcolor}
\usepackage{titlesec}
\usepackage{fancyhdr}
\usepackage{enumitem}
\usepackage{booktabs}
\usepackage{xurl}
\usepackage{lastpage}
\usepackage{etoolbox}

\definecolor{KaleidoNavy}{HTML}{0A2342}
\definecolor{KaleidoBlue}{HTML}{176B87}
\definecolor{KaleidoCyan}{HTML}{3BB7C4}
\definecolor{KaleidoMint}{HTML}{DDF5F1}
\definecolor{KaleidoFog}{HTML}{F2F5F7}
\definecolor{KaleidoSlate}{HTML}{5D6A7A}
\definecolor{KaleidoCoral}{HTML}{F26B5B}

\titleformat{\section}
  {\Large\bfseries\color{KaleidoNavy}}
  {\thesection}{0.75em}{}[\vspace{0.2em}\color{KaleidoCyan}\titlerule]
\titleformat{\subsection}
  {\large\bfseries\color{KaleidoBlue}}
  {\thesubsection}{0.65em}{}
\titleformat{\subsubsection}
  {\normalsize\bfseries\color{KaleidoNavy}}
  {\thesubsubsection}{0.55em}{}
\titlespacing*{\section}{0pt}{2.2em}{1em}
\titlespacing*{\subsection}{0pt}{1.6em}{0.65em}

\pagestyle{fancy}
\fancyhf{}
\fancyhead[L]{\small\bfseries\color{KaleidoNavy}KALEIDO FLOWTWIN}
\fancyhead[R]{\small\color{KaleidoSlate}DOSSIER TECNICO}
\fancyfoot[L]{\scriptsize\color{KaleidoSlate}EVOCON Solutions GmbH}
\fancyfoot[R]{\scriptsize\color{KaleidoSlate}\thepage/\pageref*{LastPage}}
\renewcommand{\headrulewidth}{0.6pt}
\renewcommand{\headrule}{\hbox to\headwidth{\color{KaleidoCyan}\leaders\hrule height \headrulewidth\hfill}}
\renewcommand{\footrulewidth}{0pt}
\setlength{\headheight}{20pt}

\setlist[itemize]{leftmargin=1.4em,itemsep=0.25em,topsep=0.35em}
\setlist[enumerate]{leftmargin=1.7em,itemsep=0.25em,topsep=0.35em}
\setlength{\parindent}{0pt}
\setlength{\parskip}{0.55em}
\setlength{\emergencystretch}{3em}
\renewcommand{\arraystretch}{1.18}

\AtBeginEnvironment{quote}{%
  \color{KaleidoNavy}\itshape
  \setlength{\leftskip}{1em}
  \setlength{\rightskip}{1em}
}

\makeatletter
\renewcommand{\maketitle}{%
  \hypersetup{pageanchor=false}
  \begin{titlepage}
    \pagecolor{KaleidoNavy}
    \color{white}
    \vspace*{1.4cm}
    {\color{KaleidoCyan}\rule{2.4cm}{2.5pt}\par}
    \vspace{1.1cm}
    {\Huge\bfseries\@title\par}
    \vspace{0.65cm}
    {\Large\color{KaleidoMint}Kaleido FlowTwin\par}
    \vfill
    {\large Álvaro Schwiedop Souto\par}
    \vspace{0.15cm}
    {\normalsize\color{KaleidoMint}Industrial AI \& Predictive Quality Lead\par}
    \vspace{0.15cm}
    {\normalsize\color{KaleidoMint}EVOCON Solutions GmbH\par}
    \vspace{0.25cm}
    {\normalsize\color{KaleidoMint}\@date\par}
  \end{titlepage}
  \hypersetup{pageanchor=true}
  \nopagecolor
  \color{black}
}
\makeatother
\usepackage{bookmark}
\IfFileExists{xurl.sty}{\usepackage{xurl}}{} % add URL line breaks if available
\urlstyle{same}
% fallback for those not using the hyperref driver hyperxmp:
\makeatletter
\@ifundefined{xmpquote}{\newcommand{\xmpquote}[1]{#1}}{}
\makeatother
\hypersetup{
  pdftitle={Datos y datasets},
  pdflang={es},
  colorlinks=true,
  linkcolor={KaleidoBlue},
  filecolor={Maroon},
  citecolor={Blue},
  urlcolor={KaleidoBlue},
  pdfcreator={LaTeX via pandoc}}

\title{Datos y datasets}
\author{}
\date{15 julio 2026}

\begin{document}
\maketitle

{
\setcounter{tocdepth}{3}
\tableofcontents
}
\section{Datos y datasets}\label{datos-y-datasets}

\subsection{Regla principal}\label{regla-principal}

Los datasets publicos sirven para desarrollar loaders, schemas y
baselines. Solo los datos de Kaleido pueden demostrar valor para
Kaleido.

\subsection{Prioridad A - datos internos del
piloto}\label{prioridad-a---datos-internos-del-piloto}

\subsubsection{A1. Trace Port}\label{a1-trace-port}

Campos solicitados:

\begin{itemize}
\tightlist
\item
  proyecto, operacion, turno y revision de plan;
\item
  packing list/unidades de carga;
\item
  evento/estado y timestamp;
\item
  responsable/equipo/recurso por rol;
\item
  cantidad/peso/unidad;
\item
  notas, incidencia y referencia de foto;
\item
  inicio/fin/pausas;
\item
  outcome y desviacion;
\item
  export/API version y zona horaria.
\end{itemize}

Uso: caso principal de tiempo restante y riesgo de desviacion.

\subsubsection{A2. Shipping Board}\label{a2-shipping-board}

\begin{itemize}
\tightlist
\item
  expedicion/recepcion;
\item
  milestones, documentos y alertas;
\item
  transportista y ventanas;
\item
  timestamps planificados/reales;
\item
  incidencias y entrega final.
\end{itemize}

Uso: transferencia del mismo modelo de prefijos a flujos de
instalaciones.

\subsubsection{A3. Freight Intelligence}\label{a3-freight-intelligence}

\begin{itemize}
\tightlist
\item
  contenedor, viaje, buque y carrier;
\item
  milestones DCSA-like;
\item
  ETA original y revisiones;
\item
  AIS/posicion o referencias;
\item
  alertas y excepciones;
\item
  timestamps de consulta y de ocurrencia.
\end{itemize}

Uso: ETA/excepcion, no mezclar con eventos de terminal sin mapping.

\subsubsection{A4. Coste y
planificacion}\label{a4-coste-y-planificacion}

\begin{itemize}
\tightlist
\item
  plan original y revisiones;
\item
  horas/recurso;
\item
  overtime, espera, demurrage/detention cuando aplique;
\item
  SLA/penalizacion;
\item
  coste de inspeccion y falsa alarma.
\end{itemize}

Uso: convertir metricas en utilidad. Opcional para el primer modelo,
necesario para un business case cuantitativo.

\subsection{Prioridad B - contexto
externo}\label{prioridad-b---contexto-externo}

\subsubsection{DCSA Track \& Trace y Port
Call}\label{dcsa-track--trace-y-port-call}

Estandar de procesos, eventos y APIs para transporte de contenedores.
Usar como referencia semantica, no asumir que Trace Port usa exactamente
DCSA.

\begin{itemize}
\tightlist
\item
  \url{https://dcsa.org/standards/track-and-trace/standard-documentation-track-and-trace}
\item
  \url{https://developer.dcsa.org/}
\item
  \url{https://reference.dcsa.org/content/standards/industry-blueprint/v2026-q1/industry-blueprint-2026-q1}
\end{itemize}

\subsubsection{OCEL 2.0 / XES}\label{ocel-20--xes}

OCEL es preferible para proyecto-operacion-turno-carga-recurso. XES se
puede exportar para herramientas clasicas de process mining.

\begin{itemize}
\tightlist
\item
  \url{https://www.ocel-standard.org/}
\item
  \url{https://www.tf-pm.org/resources/xes-standard/about-xes}
\end{itemize}

\subsubsection{Meteo y oceanografia}\label{meteo-y-oceanografia}

\begin{itemize}
\tightlist
\item
  Puertos del Estado/Portuscopia: oleaje, nivel, corrientes,
  temperatura, mareografos y meteo portuaria.
\item
  AEMET OpenData: observaciones/predicciones meteorologicas.
\end{itemize}

Fuentes:

\begin{itemize}
\tightlist
\item
  \url{https://portuscopia.puertos.es/}
\item
  \url{https://www.puertos.es/en/services/oceanography}
\item
  \url{https://www.aemet.es/en/datos_abiertos/AEMET_OpenData}
\end{itemize}

Verificar licencia, estacion, latencia, revision y timestamp de
publicacion. Una prediccion historica solo puede usar el forecast
disponible en ese momento, no la observacion futura.

\subsubsection{AIS}\label{ais}

NOAA MarineCadastre publica AIS de EE. UU. sin restricciones de acceso
para planificacion costera. Es util para desarrollar pipelines de
trayectoria y port-call context, no representa Vigo ni prueba el caso
Kaleido.

\begin{itemize}
\tightlist
\item
  \url{https://www.fisheries.noaa.gov/inport/item/77594}
\item
  \url{https://marinecadastre.gov/ais/}
\end{itemize}

Para Vigo se necesita una fuente legal/contractual disponible para
Kaleido.

\subsection{Prioridad C - datasets publicos para
desarrollo}\label{prioridad-c---datasets-publicos-para-desarrollo}

\subsubsection{Container Logistics Object-centric Event Log
(2026)}\label{container-logistics-object-centric-event-log-2026}

Log artificial OCEL 2.0 de coordinacion de pedido, documento, recogida,
terminal y envio. Incluye modelo CPN y SQLite/XML.

\begin{itemize}
\tightlist
\item
  Uso: tests de object graph, import OCEL, process discovery y CI.
\item
  No uso: metricas comerciales o afirmacion de generalizacion real.
\item
  Fuente: \url{https://zenodo.org/records/18373888}.
\end{itemize}

\subsubsection{BPI Challenge 2019}\label{bpi-challenge-2019}

Log real de purchase-to-pay con 251.734 items, 1.595.923 eventos y 42
actividades, anonimizado y XES.

\begin{itemize}
\tightlist
\item
  Uso: escalabilidad, remaining time/next-event y baselines PPM.
\item
  No uso: evidencia portuaria.
\item
  DOI:
  \url{https://doi.org/10.4121/uuid:d06aff4b-79f0-45e6-8ec8-e19730c248f1}.
\end{itemize}

\subsubsection{Outbound Warehouse
Processes}\label{outbound-warehouse-processes}

Caso de 2025 con 169.523 trazas de un proceso real de almacen outbound.
El paper indica disponibilidad publica; antes de descargar se debe
localizar el release oficial, version y licencia.

\begin{itemize}
\tightlist
\item
  Uso: comparacion boosting vs deep PPM.
\item
  Fuente primaria: \url{https://arxiv.org/abs/2509.18986}.
\end{itemize}

\subsubsection{LaDe}\label{lade}

Dataset industrial de last-mile con paquetes y task events multi-ciudad.

\begin{itemize}
\tightlist
\item
  Uso: generalizacion espaciotemporal/last-mile opcional.
\item
  No uso: terminal portuaria.
\item
  Fuente: \url{https://arxiv.org/abs/2306.10675}.
\end{itemize}

\subsubsection{BPI logs adicionales}\label{bpi-logs-adicionales}

Solo para pretraining/benchmark de event encoders. Separar por dataset y
no presentar un promedio que oculte fallos por log.

\begin{itemize}
\tightlist
\item
  \url{https://www.tf-pm.org/resources/logs}
\end{itemize}

\subsection{Datos sinteticos}\label{datos-sinteticos}

La simulacion discreta y el CPN pueden producir eventos para:

\begin{itemize}
\tightlist
\item
  tests y fixtures;
\item
  fallos raros controlados;
\item
  comprobar invariantes y escenarios;
\item
  pretraining experimental claramente etiquetado.
\end{itemize}

No se usan para estimar precision real, ROI ni frecuencia de
incidencias.

\subsection{Manifest obligatorio}\label{manifest-obligatorio}

\begin{Shaded}
\begin{Highlighting}[]
\FunctionTok{dataset\_id}\KeywordTok{:}\AttributeTok{ trace\_port\_export\_v1}
\FunctionTok{owner}\KeywordTok{:}\AttributeTok{ kaleido}
\FunctionTok{source\_system}\KeywordTok{:}\AttributeTok{ trace\_port}
\FunctionTok{export\_version}\KeywordTok{:}\AttributeTok{ unknown}
\FunctionTok{access\_date}\KeywordTok{:}\AttributeTok{ YYYY{-}MM{-}DD}
\FunctionTok{license\_or\_agreement}\KeywordTok{:}\AttributeTok{ pending}
\FunctionTok{timezone\_source}\KeywordTok{:}\AttributeTok{ Europe/Madrid}
\FunctionTok{rows}\KeywordTok{:}\AttributeTok{ measured}
\FunctionTok{operations}\KeywordTok{:}\AttributeTok{ measured}
\FunctionTok{projects}\KeywordTok{:}\AttributeTok{ measured}
\FunctionTok{date\_min}\KeywordTok{:}\AttributeTok{ measured}
\FunctionTok{date\_max}\KeywordTok{:}\AttributeTok{ measured}
\FunctionTok{sha256}\KeywordTok{:}\AttributeTok{ measured}
\FunctionTok{contains\_personal\_data}\KeywordTok{:}\AttributeTok{ unknown}
\FunctionTok{contains\_photos}\KeywordTok{:}\AttributeTok{ }\CharTok{false}
\FunctionTok{plan\_revisions\_available}\KeywordTok{:}\AttributeTok{ unknown}
\FunctionTok{outcomes\_available}\KeywordTok{:}\AttributeTok{ unknown}
\FunctionTok{action\_columns}\KeywordTok{:}\AttributeTok{ }\KeywordTok{[]}
\FunctionTok{context\_columns}\KeywordTok{:}\AttributeTok{ }\KeywordTok{[]}
\FunctionTok{observation\_columns}\KeywordTok{:}\AttributeTok{ }\KeywordTok{[]}
\FunctionTok{outcome\_columns}\KeywordTok{:}\AttributeTok{ }\KeywordTok{[]}
\FunctionTok{forbidden\_columns}\KeywordTok{:}\AttributeTok{ }\KeywordTok{[]}
\FunctionTok{known\_limitations}\KeywordTok{:}\AttributeTok{ }\KeywordTok{[]}
\end{Highlighting}
\end{Shaded}

\subsection{Quality gates}\label{quality-gates}

\begin{itemize}
\tightlist
\item
  IDs estables y cobertura de relaciones;
\item
  timezone y orden temporal;
\item
  duplicados/reintentos;
\item
  plan vs actual distinguibles;
\item
  granularidad/event vocabulary;
\item
  missingness por proyecto y periodo;
\item
  outcome y censura;
\item
  accion/context audit;
\item
  riesgo de leakage;
\item
  representatividad por carga/cliente/puerto;
\item
  consentimiento/licencia/retencion.
\end{itemize}

\subsection{Presupuesto de datos para
go/no-go}\label{presupuesto-de-datos-para-gono-go}

No se exige un numero magico. La decision usa:

\begin{itemize}
\tightlist
\item
  operaciones completas;
\item
  diversidad de variantes;
\item
  numero de eventos por prefijo;
\item
  cobertura de plan y outcome;
\item
  suficientes desviaciones/incidencias;
\item
  estabilidad de IDs;
\item
  separacion temporal disponible.
\end{itemize}

Si faltan outcomes, el piloto pasa a \texttt{instrument-first}: process
mining, simulacion y plan de captura, sin clasificador de riesgo.

\end{document}
